\documentclass[10pt]{article}
\usepackage[
  paperwidth=8.5in,
  paperheight=11in,
  top=0.85in,
  bottom=1in,
  left=2.5in,
  right=0.85in,
  marginparwidth=1.75in,
  marginparsep=0.25in,
  headheight=14pt
]{geometry}
\reversemarginpar

\usepackage{amsmath,amssymb,amsfonts}
\usepackage{graphicx}
\usepackage{xcolor}
\usepackage{fontspec}
\usepackage{fancyhdr}
\usepackage{marginnote}

\IfFontExistsTF{Times New Roman}{
  \setmainfont{Times New Roman}
}{
  \IfFontExistsTF{TeX Gyre Termes}{
    \setmainfont{TeX Gyre Termes}
  }{}
}
\IfFontExistsTF{Arial}{
  \setsansfont{Arial}
}{
  \IfFontExistsTF{TeX Gyre Heros}{
    \setsansfont{TeX Gyre Heros}
  }{}
}

\definecolor{stewartblue}{RGB}{0, 130, 195}
\definecolor{stewartred}{RGB}{204, 0, 0}

\pagestyle{fancy}
\fancyhf{}
\renewcommand{\headrulewidth}{0pt}
\renewcommand{\footrulewidth}{0pt}
\fancyhead[R]{\textsf{\small MỤC 2.5 \quad Tính liên tục} \qquad \textbf{\textsf{119}}}

\fancyfootoffset[L]{\dimexpr\marginparwidth+\marginparsep\relax}
\fancyfoot[C]{%
  \begin{minipage}{\dimexpr\textwidth+\marginparwidth+\marginparsep\relax}
    \tiny\color{black!50}\centering
    Copyright 2021 Cengage Learning. All Rights Reserved. May not be copied, scanned, or duplicated, in whole or in part. WCN 02-200-203\\[1pt]
    Copyright 2021 Cengage Learning. All Rights Reserved. May not be copied, scanned, or duplicated, in whole or in part. Due to electronic rights, some third party content may be suppressed from the eBook and/or eChapter(s). Editorial review has deemed that any suppressed content does not materially affect the overall learning experience. Cengage Learning reserves the right to remove additional content at any time if subsequent rights restrictions require it.
  \end{minipage}%
}

\setlength{\parindent}{15pt}
\setlength{\parskip}{0pt}

\begin{document}

\noindent
\textbf{(b)}\enskip Một hàm phân thức hữu tỉ là một hàm số có dạng
\[
f(x) = \frac{P(x)}{Q(x)}
\]
trong đó $P$ và $Q$ là các đa thức. Tập xác định của $f$ là $D = \{x \in \mathbb{R} \mid Q(x) \neq 0\}$. Ta đã biết từ phần (a) rằng $P$ và $Q$ liên tục khắp nơi. Do đó, theo phần 5 của Định lý 4, $f$ liên tục tại mọi số thuộc $D$.\hfill\textcolor{stewartred}{\rule{1.1ex}{1.1ex}}

\medskip
Để minh họa cho Định lý 5, hãy quan sát thấy rằng thể tích của một hình cầu biến thiên liên tục theo bán kính của nó bởi vì công thức $V(r) = \frac{4}{3}\pi r^3$ cho thấy $V$ là một hàm đa thức của $r$. Tương tự, nếu một quả bóng được ném thẳng đứng lên không trung với vận tốc ban đầu là $50\text{ ft/s}$, thì độ cao của quả bóng tính bằng feet sau $t$ giây được cho bởi công thức $h = 50t - 16t^2$. Một lần nữa, đây là một hàm đa thức, vì vậy độ cao là một hàm liên tục của thời gian trôi qua, đúng như chúng ta kỳ vọng.

Việc biết được những hàm số nào liên tục cho phép chúng ta tính một số giới hạn rất nhanh chóng, như ví dụ sau đây sẽ chỉ ra. Hãy so sánh ví dụ này với Ví dụ 2.3.2(b).

\bigskip
\noindent
{\color{stewartblue}\textbf{\textsf{VÍ DỤ 5}}}\enskip Tìm $\lim\limits_{x \to -2} \dfrac{x^3 + 2x^2 - 1}{5 - 3x}$.

\medskip\noindent
{\color{stewartblue}\textbf{\textsf{LỜI GIẢI}}}\enskip Hàm số
\[
f(x) = \frac{x^3 + 2x^2 - 1}{5 - 3x}
\]
là hàm phân thức hữu tỉ, nên theo Định lý 5 nó liên tục trên tập xác định của nó, tức là $\{x \mid x \neq \frac{5}{3}\}$. Do đó
\begin{align*}
\lim_{x \to -2} \frac{x^3 + 2x^2 - 1}{5 - 3x} &= \lim_{x \to -2} f(x) = f(-2) \\[4pt]
&= \frac{(-2)^3 + 2(-2)^2 - 1}{5 - 3(-2)} = -\frac{1}{11} \tag*{{\color{stewartblue}\blacksquare}}
\end{align*}
\marginnote{%
  \centering
  \includegraphics[width=\linewidth]{images/p0154_fig1.png}\\[3pt]
  {\small\textbf{\textsf{HÌNH 5}}}
}[-1.0cm]%
Hóa ra hầu hết các hàm số quen thuộc đều liên tục tại mọi số thuộc tập xác định của chúng. Chẳng hạn, Quy tắc Giới hạn 11 trong Mục 2.3 chính là khẳng định rằng các hàm căn thức là liên tục.

Từ hình dạng đồ thị của các hàm số sin và cosin (Hình 1.2.19), chúng ta chắc chắn sẽ đoán rằng chúng liên tục. Ta biết từ định nghĩa của $\sin\theta$ và $\cos\theta$ rằng tọa độ của điểm $P$ trong Hình 5 là $(\cos\theta, \sin\theta)$. Khi $\theta \to 0$, ta thấy điểm $P$ tiến dần về điểm $(1, 0)$ và vì vậy $\cos\theta \to 1$ và $\sin\theta \to 0$. Do đó

\begin{equation*}
\raisebox{0.1ex}{\fboxrule=0.8pt\fboxsep=2.5pt\fcolorbox{stewartred}{white}{\textcolor{stewartred}{\textbf{6}}}}
\hfill
\lim_{\theta \to 0} \cos\theta = 1 \qquad\qquad \lim_{\theta \to 0} \sin\theta = 0
\hfill
\phantom{\raisebox{0.1ex}{\fboxrule=0.8pt\fboxsep=2.5pt\fcolorbox{stewartred}{white}{\textbf{6}}}}
\end{equation*}

\noindent
\marginnote{%
  \raggedright\footnotesize
  Một cách khác để thiết lập các giới hạn trong (6) là sử dụng Định lý Kẹp với bất đẳng thức $\sin\theta < \theta$ (với $\theta > 0$), điều này được chứng minh trong Mục 3.3.
}[0pt]%
Vì $\cos 0 = 1$ và $\sin 0 = 0$, các phương trình trong (6) khẳng định rằng các hàm cosin và sin liên tục tại 0. Sau đó, các công thức cộng cho cosin và sin có thể được dùng để suy ra rằng các hàm này liên tục khắp nơi (xem Bài tập 66 và 67).

Từ phần 5 của Định lý 4, ta suy ra rằng
\[
\tan x = \frac{\sin x}{\cos x}
\]
\noindent
liên tục ngoại trừ tại những điểm mà $\cos x = 0$. Điều này xảy ra khi $x$ là một bội số nguyên lẻ

\end{document}